% ID: SEL-D0054-OB-E02
% ESTADO: cotejada. Segunda pasada visual y borrador compilado revisados.
% Fuente: Selectividad/sel_2011_matematicas_aplicadas/121-EXAMEN 4.pdf
% SHA-256: c70ed065da193c82dd0a02fb4c05069415b0adf53d0f49bba1b0375a19823566
% Página: 2; opción B; ejercicio original 2.

Las funciones $I(t)=-2t^2+51t$ y $G(t)=t^2-3t+96$ con $0\leq t\leq18$ representan, respectivamente, los ingresos y gastos de una empresa, en miles de euros, en función de los años, $t$, transcurridos desde su inicio y en los últimos 18 años.
\begin{enumerate}[label=\alph*)]
\item \textbf{(0.5 puntos)} \textquestiondown{}Para qué valores de $t$, desde su entrada en funcionamiento, los ingresos coincidieron con los gastos?
\item \textbf{(1 punto)} Determine la función que refleje los beneficios (ingresos menos gastos) en función de $t$ y represéntela gráficamente.
\item \textbf{(1 punto)} \textquestiondown{}Al cabo de cuántos años, desde su entrada en funcionamiento, los beneficios fueron máximos? Calcule el valor de ese beneficio.
\end{enumerate}
